\section*{Chapitre \ref{ch:g11:cell-cycle-mitosis} --- Le cycle cellulaire et la mitose}

\begin{solution}{exo:g11:cell-cycle-mitosis:1}
$\mathrm{G_1}$~: croissance~; S~: \omterm{def:g11:dna-replication:replication}{réplication de l'ADN}~;
$\mathrm{G_2}$~: préparation de la division~; M~: \omterm{def:g11:cell-cycle-mitosis:cycle}{mitose} (division du
\omterm{def:g10:cells-common-unit:organelle}{noyau}) et cytodiérèse (division du \omterm{def:g10:cells-common-unit:cell}{cytoplasme}).
\end{solution}

\begin{solution}{exo:g11:cell-cycle-mitosis:2}
Une \omterm{def:g11:cell-cycle-mitosis:chromosome}{chromatide} est l'une des deux copies identiques d'un \omterm{def:g11:cell-cycle-mitosis:chromosome}{chromosome}
répliqué~; le \omterm{def:g11:cell-cycle-mitosis:chromosome}{centromère} est la région où les deux sont maintenues
ensemble. Un \omterm{def:g11:cell-cycle-mitosis:chromosome}{chromosome} a deux \omterm{def:g11:cell-cycle-mitosis:chromosome}{chromatides} de la fin de la phase S
jusqu'à l'anaphase.
\end{solution}

\begin{solution}{exo:g11:cell-cycle-mitosis:3}
Prophase~: les \omterm{def:g10:universal-dna:information}{chromosomes} se condensent, l'enveloppe nucléaire se
rompt. Métaphase~: les \omterm{def:g10:universal-dna:information}{chromosomes} s'alignent sur l'équateur.
Anaphase~: les \omterm{def:g11:cell-cycle-mitosis:chromosome}{chromatides} se séparent et gagnent les pôles.
Télophase~: les enveloppes nucléaires se reforment, puis vient la
cytodiérèse.
\end{solution}

\begin{solution}{exo:g11:cell-cycle-mitosis:4}
En phase S, à mi-chemin de la réplication.
\end{solution}

\begin{solution}{exo:g11:cell-cycle-mitosis:5}
46 \omterm{def:g10:universal-dna:information}{chromosomes}, 92 \omterm{def:g11:cell-cycle-mitosis:chromosome}{chromatides}.
\end{solution}

\begin{solution}{exo:g11:cell-cycle-mitosis:6}
$\mathrm{G_1}$ \qty{11}{h}, S \qty{8}{h}, $\mathrm{G_2}$ \qty{4}{h},
M \qty{1}{h}. Une \omterm{def:g10:cells-common-unit:cell}{cellule} à $2q$ est en $\mathrm{G_2}$ ou en \omterm{def:g11:cell-cycle-mitosis:cycle}{mitose}
avant la fin de l'anaphase.
\end{solution}

\begin{solution}{exo:g11:cell-cycle-mitosis:7}
Métaphase~: 40 \omterm{def:g10:universal-dna:information}{chromosomes}, 80 \omterm{def:g11:cell-cycle-mitosis:chromosome}{chromatides}. Anaphase~: 80 \omterm{def:g10:universal-dna:information}{chromosomes}
(à une \omterm{def:g11:cell-cycle-mitosis:chromosome}{chromatide}), 40 vers chaque pôle. \omterm{def:g10:cells-common-unit:cell}{Cellule} fille~: 40
\omterm{def:g10:universal-dna:information}{chromosomes}, 40 \omterm{def:g11:cell-cycle-mitosis:chromosome}{chromatides}.
\end{solution}

\begin{solution}{exo:g11:cell-cycle-mitosis:8}
Index mitotique $80/800 = 10\%$~; \omterm{def:g11:cell-cycle-mitosis:cycle}{mitose} \qty{2.4}{h}. Prophase
$48/800 \times 24 = \qty{1.44}{h}$~; métaphase \qty{0.48}{h}
(29 min)~; anaphase \qty{0.18}{h} (11 min)~; télophase \qty{0.3}{h}
(18 min).
\end{solution}

\begin{solution}{exo:g11:cell-cycle-mitosis:9}
Sans fuseau, les \omterm{def:g11:cell-cycle-mitosis:chromosome}{chromatides} ne peuvent pas être écartées~: la \omterm{def:g10:cells-common-unit:cell}{cellule}
reste bloquée en métaphase, ses \omterm{def:g10:universal-dna:information}{chromosomes} condensés et parfaitement
visibles, chacun avec ses deux \omterm{def:g11:cell-cycle-mitosis:chromosome}{chromatides}. C'est exactement l'état
dans lequel on photographie les \omterm{def:g10:universal-dna:information}{chromosomes} pour un \omterm{def:g11:cell-cycle-mitosis:chromosome}{caryotype}.
\end{solution}

\begin{solution}{exo:g11:cell-cycle-mitosis:10}
Deux mètres de fil emmêlé ne peuvent pas être triés~; condensés en
bâtonnets de quelques micromètres, les 46 \omterm{def:g10:universal-dna:information}{chromosomes} peuvent être
alignés et écartés sans se rompre ni se nouer. Ils se déroulent ensuite
parce que les \omterm{def:g10:universal-dna:gene}{gènes} ne peuvent être lus que sous la forme relâchée.
\end{solution}

\begin{solution}{exo:g11:cell-cycle-mitosis:11}
Dix divisions~: $2^{10} = 1024$ \omterm{def:g10:cells-common-unit:cell}{cellules}. Non~: l'embryon précoce ne
croît pas entre deux divisions, et ses cycles sont presque uniquement
faits de S et de M, avec quasiment pas de $\mathrm{G_1}$.
\end{solution}

\begin{solution}{exo:g11:cell-cycle-mitosis:12}
Une fille a 47 \omterm{def:g10:universal-dna:information}{chromosomes} (un de trop), l'autre 45 (un de moins)~;
leur \omterm{def:g10:universal-dna:information}{ADN} vaut $q$ plus ou moins la valeur d'un \omterm{def:g11:cell-cycle-mitosis:chromosome}{chromosome}. Dans une
\omterm{def:g10:cells-common-unit:cell}{cellule} du corps, l'erreur ne touche qu'une \omterm{def:g10:cells-common-unit:cell}{cellule} parmi des milliers
de milliards, qui d'ordinaire meurt ou est éliminée~; mais si la
\omterm{def:g10:cells-common-unit:cell}{cellule} survit et se divise, et que le \omterm{def:g11:cell-cycle-mitosis:chromosome}{chromosome} perdu ou gagné
dérègle sa croissance, une tumeur peut démarrer.
\end{solution}

\begin{solution}{exo:g11:cell-cycle-mitosis:13}
Environ $\num{2e11}$ divisions par jour, soit $\num{2e11}$ \omterm{def:g10:cells-common-unit:cell}{cellules}
dans le cycle à un instant donné (chacune rendant une \omterm{def:g10:cells-common-unit:cell}{cellule} neuve par
jour). Bloquer la \omterm{def:g11:cell-cycle-mitosis:cycle}{mitose} coupe l'approvisionnement~; les globules
rouges vivent 120 jours, mais certains globules blancs quelques jours
seulement, si bien qu'en une semaine le nombre de globules blancs
s'effondre et que les infections suivent --- l'effet secondaire
classique des médicaments anticancéreux.
\end{solution}

\begin{solution}{exo:g11:cell-cycle-mitosis:14}
En interphase, les \omterm{def:g10:universal-dna:information}{chromosomes} sont déroulés et invisibles~; en
anaphase, ils sont en mouvement et mêlés. En métaphase, ils sont
entièrement condensés, séparés, encore faits de deux \omterm{def:g11:cell-cycle-mitosis:chromosome}{chromatides} et
disposés dans un même plan~: chacun peut être identifié par sa taille
et sa forme.
\end{solution}

\begin{solution}{exo:g11:cell-cycle-mitosis:15}
Des \omterm{def:g10:cells-common-unit:cell}{cellules} hépatiques au repos en $\mathrm{G_1}$ sont rentrées dans
le cycle, ont traversé S, $\mathrm{G_2}$ et la \omterm{def:g11:cell-cycle-mitosis:cycle}{mitose} à plusieurs
reprises jusqu'à ce que la masse soit rétablie, puis sont retournées au
repos. Au jour 0, presque toutes les \omterm{def:g10:cells-common-unit:cell}{cellules} détenaient $q$~; au jour
3, beaucoup détenaient entre $q$ et $2q$ (en S) ou $2q$ (en
$\mathrm{G_2}$, en \omterm{def:g11:cell-cycle-mitosis:cycle}{mitose}).
\end{solution}

\begin{solution}{pb:g11:cell-cycle-mitosis:1}
\textbf{1.} $100/1000 = 10\%$.

\textbf{2.} $0.10 \times 20 = \qty{2}{h}$.

\textbf{3.} Prophase $60/1000 \times 20 = \qty{1.2}{h}$ (72 min)~;
métaphase $0.36$ h (22 min)~; anaphase $0.16$ h (environ 10 min)~;
télophase $0.28$ h (17 min).

\textbf{4.} L'anaphase. Les \omterm{def:g11:cell-cycle-mitosis:chromosome}{chromatides} sont déjà condensées et
attachées~; le mouvement est une traction de quelques micromètres par
des fibres qui se raccourcissent en quelques minutes.

\textbf{5.} Les stades rares ne sont vus que dans peu de \omterm{def:g10:cells-common-unit:cell}{cellules}, et
un petit effectif donne une grande erreur~; et seule la zone de
croissance boucle des cycles --- ailleurs l'index serait nul et
diluerait les chiffres.

\textbf{6.} 16 \omterm{def:g10:universal-dna:information}{chromosomes}, 32 \omterm{def:g11:cell-cycle-mitosis:chromosome}{chromatides}.

\textbf{7.} 32 \omterm{def:g10:universal-dna:information}{chromosomes} en mouvement, 16 vers chaque pôle.

\textbf{8.} Prophase $2q$~; anaphase $2q$ pour la \omterm{def:g10:cells-common-unit:cell}{cellule} entière~;
télophase $q$ par \omterm{def:g10:cells-common-unit:organelle}{noyau}.

\textbf{9.} S occupe $7/19$ de l'interphase~: environ $900 \times 7/19 \approx
330$ \omterm{def:g10:cells-common-unit:cell}{cellules}.

\textbf{10.} L'anaphase~: les \omterm{def:g11:cell-cycle-mitosis:chromosome}{centromères} viennent de se cliver et les
32 \omterm{def:g11:cell-cycle-mitosis:chromosome}{chromatides} des 16 \omterm{def:g10:universal-dna:information}{chromosomes} sont devenues 32 \omterm{def:g10:universal-dna:information}{chromosomes} en route
vers les pôles.

\textbf{11.} $\num{20000}/20 = 1000$ \omterm{def:g10:cells-common-unit:cell}{cellules} neuves par heure.

\textbf{12.} 1000 \omterm{def:g10:cells-common-unit:cell}{cellules} par heure sur toute la zone~; en file, la
zone fait peut-être une centaine de \omterm{def:g10:cells-common-unit:cell}{cellules} de large dans chaque
direction, soit à peu près
$\num{24000}/(100 \times 100) \approx 2$--3 \omterm{def:g10:cells-common-unit:cell}{cellules} par file et par
jour, de \qty{100}{\micro m} chacune~: quelques dixièmes de
millimètre. Le \qty{1}{mm} observé signifie que la zone est plus
étroite qu'on ne l'a supposé, ou que les \omterm{def:g10:cells-common-unit:cell}{cellules} s'allongent
davantage~; l'ordre de grandeur est le bon.

\textbf{13.} Une origine copie $2 \times 50 \times 7 \times 3600 =
\num{2.5e6}$ paires de \omterm{def:g10:universal-dna:nucleotide}{bases}~; $\num{1.6e10}/\num{2.5e6} \approx 6400$
origines au minimum.

\textbf{14.} Chaque stade dure plus longtemps, de sorte que moins de
\omterm{def:g10:cells-common-unit:cell}{cellules} bouclent le cycle par jour~; si toutes les phases ralentissent
également, les proportions, et donc l'index mitotique, ne changent pas.

\textbf{15.} Les \omterm{def:g10:cells-common-unit:cell}{cellules} déjà sorties de S achèveraient $\mathrm{G_2}$
et la \omterm{def:g11:cell-cycle-mitosis:cycle}{mitose} dans les premières heures~; au bout d'un jour, aucune
\omterm{def:g10:cells-common-unit:cell}{cellule} ne peut plus atteindre la \omterm{def:g11:cell-cycle-mitosis:cycle}{mitose}, si bien que de la prophase à
la télophase tout a disparu et qu'il ne reste que des \omterm{def:g10:cells-common-unit:cell}{cellules} en
interphase, bloquées en $\mathrm{G_1}$ ou au début de S.

\textbf{16.} $4/1000 \times 20 = \qty{0.08}{h}$, soit environ 5
minutes, contre 10. Rien de surprenant~: avec 4 ou 8 \omterm{def:g10:cells-common-unit:cell}{cellules}, le
comptage est dominé par le hasard.

\textbf{17.} $12/2000 \times 20 = \qty{0.12}{h}$, soit environ 7 minutes.

\textbf{18.} Humain~: $0.04 \times 24 \approx \qty{1}{h}$~; ail
\qty{2}{h}. Ordre de grandeur voisin, les \omterm{def:g10:cells-common-unit:cell}{cellules} végétales étant un
peu plus lentes.

\textbf{19.} Les \omterm{def:g10:cells-common-unit:cell}{cellules} tumorales bouclent des cycles en continu
alors que la plupart des \omterm{def:g10:cells-common-unit:cell}{cellules} normales se reposent en
$\mathrm{G_1}$, si bien qu'une fraction plus grande d'entre elles est
en \omterm{def:g11:cell-cycle-mitosis:cycle}{mitose} à un instant donné. Un médicament qui bloque le fuseau touche
donc les \omterm{def:g10:cells-common-unit:cell}{cellules} tumorales plus souvent que les normales --- mais
aussi les tissus sains qui se divisent vite.

\textbf{20.} Cycle de 20 heures~: $\mathrm{G_1}$ 8, S 7, $\mathrm{G_2}$
4, \omterm{def:g11:cell-cycle-mitosis:cycle}{mitose} environ 2 (prophase 72 min, métaphase 22, anaphase 7--10,
télophase 17). L'anaphase, la plus courte, est le moment où les
\omterm{def:g11:cell-cycle-mitosis:chromosome}{centromères} se clivent et où chaque pôle reçoit une \omterm{def:g11:cell-cycle-mitosis:chromosome}{chromatide} de
chacun des seize \omterm{def:g10:universal-dna:information}{chromosomes}.
\end{solution}
